\documentclass[12pt]{article}

\usepackage[margin=1in]{geometry}
\usepackage{amsmath,amsthm,amssymb}
\usepackage{fancyhdr}
\usepackage{graphicx}
\usepackage{enumitem}
\usepackage{bm}
\pagestyle{fancy}
\setlength{\headheight}{28pt}

\usepackage{xcolor}
\usepackage[most]{tcolorbox}
\usepackage{xparse}
\tcbuselibrary{skins,theorems}

% Macro Bible-style remark environment (boxed)
\tcbset{
  boxbase/.style={
    enhanced,breakable,
    fonttitle=\bfseries,
    boxrule=.5pt,arc=2pt,boxsep=5pt,
    before skip=10pt,after skip=10pt},
  white/.style={
    colback=white,
    colframe=black,
    coltitle=white,
    colbacktitle=black}
}
\newtcbtheorem[auto counter]{remark}{Remark}{boxbase,white}{rem}
\let\TCBremark      \remark
\let\endTCBremark   \endremark
\renewenvironment{remark}[1][]
  {\IfValueTF{#1}
     {\TCBremark[]{#1}{}}
     {\TCBremark[][]{}}}
  {\endTCBremark}

% Common shortcuts
\newcommand{\R}{\mathbb R}
\newcommand{\N}{\mathbb N}
\newcommand{\Z}{\mathbb Z}
\newcommand{\E}{\mathbb E}
\newcommand{\Var}{\mathrm{Var}}
\newcommand{\Cov}{\mathrm{Cov}}
\newcommand{\tr}{\mathrm{tr}}
\newcommand{\rank}{\mathrm{rank}}
\newcommand{\vect}{\mathrm{vec}}
\newcommand{\diag}{\mathrm{diag}}
\DeclareMathOperator*{\plim}{plim}
\DeclareMathOperator*{\argmin}{arg\,min}
\DeclareMathOperator*{\argmax}{arg\,max}

\usepackage{datetime}
\newdate{date}{27}{03}{2026}

\usepackage[round]{natbib}
\usepackage[colorlinks=true, pdfstartview=FitV, linkcolor=blue,
citecolor=blue, urlcolor=blue]{hyperref}

\theoremstyle{definition}
\newtheorem{solution}{Solution to Problem}

\begin{document}

\lhead{TA: Rafael Lincoln \\ Prof. Tim Cogley}
\rhead{Econometrics II (Time Series) \\ Solutions: \displaydate{date}}

\section*{Problem Set 5\footnote{These solutions rely partly on the past solutions of Bálint Szőke for Tim Christensen's ECON GA 2101:Econometrics II course} --- Solutions}
\subsection*{GMM: Identification, Testing, Misspecification, and Optimal Instruments}

\vspace{1em}

%=============================================================================
% SOLUTION TO PROBLEM 1
%=============================================================================
\begin{solution}[Testing a subset of moment conditions: Eichenbaum--Hansen--Singleton (1988)] \citet{EichenbaumHansenSingleton1988} estimate a representative-agent model of consumption and leisure choice. The agent chooses consumption and leisure streams $(c_t,\ell_t)_{t\geq 0}$ to maximize the time- and state-separable utility function
\[
\E\left[ \sum_{t=0}^{\infty} \beta^t u(c_t,\ell_t;\gamma,\alpha) \right]
\]
subject to a budget constraint. This leads to first-order conditions that can be written as
\[
\E\left[ h_1\!\left( \frac{c_{t+1}}{c_t}, \frac{\ell_t}{c_t}; \beta, \gamma, \alpha \right) R_{t+1} \mid \mathcal{F}_t \right] = 1, \qquad
\E\left[ h_2\!\left( \frac{c_{t+1}}{c_t}, \frac{\ell_t}{c_t}; \beta, \gamma, \alpha \right) R_{t+1} \mid \mathcal{F}_t \right] = 1,
\]
where $h_1$ and $h_2$ are known functions, $(\beta,\gamma,\alpha)$ are unknown preference parameters in some parameter space $\Theta\subset\R^3$, and $R_{t+1}$ is the return on some asset. The authors estimate the true preference parameters by GMM using the population moment conditions
\[
\E\left[ \begin{pmatrix} h_1(\cdot;\beta,\gamma,\alpha) R_{t+1} - 1 \\ h_2(\cdot;\beta,\gamma,\alpha) R_{t+1} - 1 \end{pmatrix} \otimes z_t \right] = 0,
\]
where $z_t$ is a $4\times 1$ vector of variables in $\mathcal{F}_t$. You may assume that the data $w_1,\ldots,w_n$ are strictly stationary and ergodic (with $w_t$ containing $c_{t+1}/c_t$, $\ell_t/c_t$, $R_{t+1}$, and $z_t$), and that all usual regularity conditions hold.

The presence of labor market frictions may make the moment conditions involving $h_2$ invalid, i.e.\ it may be that
\[
\E\left[ \bigl( h_2(\cdot;\beta_0,\gamma_0,\alpha_0) R_{t+1} - 1 \bigr) \otimes z_t \right] \neq 0.
\]
The question asks: describe how you would test whether the moment conditions involving $h_2$ are invalid; state (i) the null and alternative hypotheses, (ii) the test statistic, and (iii) its asymptotic distribution under the null.

\bigskip

\noindent\textbf{Solution.} First, we partition the moment conditions into two blocks:
\[
g(w_t;\theta) = \begin{pmatrix} g_1(w_t;\theta) \\ g_2(w_t;\theta) \end{pmatrix}, \qquad g_1\in\R^4,\quad g_2\in\R^4,
\]
where $g_1(w_t;\theta) = (h_1(\cdot;\theta)R_{t+1}-1)\otimes z_t$ collects the four moment conditions from the consumption Euler equation, and $g_2(w_t;\theta) = (h_2(\cdot;\theta)R_{t+1}-1)\otimes z_t$ collects the four moment conditions from the labor--leisure Euler equation. Under the null hypothesis, both blocks are valid at $\theta_0$; under the alternative, only $g_1$ is valid.

\bigskip
\paragraph{Hypotheses.}
\begin{itemize}
    \item $H_0$: All $m=8$ moment conditions are correctly specified, i.e.\ $\E[g(w_t;\theta_0)]=0$.
    \item $H_1$: The moment conditions involving $h_2$ are invalid, i.e.\ $\E[g_2(w_t;\theta_0)]\neq 0$ while $\E[g_1(w_t;\theta_0)]=0$.
\end{itemize}

\bigskip
\paragraph{Test statistic: the $C$-test (incremental $J$-test).} The idea is to compare the Hansen $J$-statistic from the full model (using all 8 moments) with the $J$-statistic from a restricted model (using only the 4 ``safe'' moments from $g_1$). If the suspect moments $g_2$ are valid, adding them should not significantly worsen the fit of the model.

\medskip
\noindent\textbf{Step 1: Full-model $J$-statistic.} Estimate $\theta$ by efficient two-step GMM using all $m=8$ moment conditions. Let $\hat{\theta}_{\text{full}}$ denote the resulting estimator, and define
\[
J_{\text{full}} = n\,\bar{g}_n(\hat{\theta}_{\text{full}})'\hat{S}_{\text{full}}^{-1}\bar{g}_n(\hat{\theta}_{\text{full}}),
\]
where $\hat{S}_{\text{full}}$ is a consistent estimator of the long-run variance $S=\Var(\sqrt{n}\,\bar{g}_n(\theta_0))$ based on the full set of moments. Under $H_0$:
\[
J_{\text{full}} \xrightarrow{d} \chi^2(m-p) = \chi^2(5).
\]

\noindent\textbf{Step 2: Restricted-model $J$-statistic.} Estimate $\theta$ by efficient two-step GMM using only the $m_1=4$ moment conditions from $g_1$. Let $\hat{\theta}_{\text{rest}}$ denote this estimator, and define
\[
J_{\text{rest}} = n\,\bar{g}_{1,n}(\hat{\theta}_{\text{rest}})'\hat{S}_{11}^{-1}\bar{g}_{1,n}(\hat{\theta}_{\text{rest}}),
\]
where $\hat{S}_{11}$ is the upper-left $4\times 4$ block of $\hat{S}_{\text{full}}$ (the long-run variance of the $g_1$ moments). Under $H_0$:
\[
J_{\text{rest}} \xrightarrow{d} \chi^2(m_1-p) = \chi^2(1).
\]

\noindent\textbf{Step 3: The $C$-statistic.} Define
\[
\boxed{C = J_{\text{full}} - J_{\text{rest}}.}
\]
Under $H_0$, both $J$-statistics are asymptotically $\chi^2$ and, by a result due to Newey (1985) and Eichenbaum, Hansen, and Singleton (1988), their difference is asymptotically independent of $J_{\text{rest}}$. The degrees of freedom of $C$ equal the number of suspect moment conditions:
\[
\boxed{C \xrightarrow{d} \chi^2(m_2) = \chi^2(4) \quad\text{under }H_0,}
\]
where $m_2=m-m_1=4$ is the number of moment conditions being tested (those involving $h_2$). We reject $H_0$ at significance level $\alpha$ if $C>\chi^2_{4,1-\alpha}$.
\end{solution}


\bigskip
\paragraph{Intuition.} The $C$-test asks: \emph{how much does the $J$-statistic worsen when we add the suspect moments?} Under the null, adding valid moments should not create a significant increase in the over-identification test statistic---any increase is due to sampling variation and is asymptotically $\chi^2(4)$. Under the alternative, the suspect moments are misspecified, so forcing the GMM estimator to fit them as well will substantially inflate the $J$-statistic, producing a large value of $C$.

\paragraph{Important implementation detail.} For the difference $C=J_{\text{full}}-J_{\text{rest}}$ to be guaranteed non-negative in finite samples (and for the asymptotic theory to apply cleanly), both $J$-statistics should be constructed using the \emph{same} consistent estimator of the long-run variance matrix $S$. In practice, one typically estimates $S$ once from the full model and uses the appropriate sub-blocks for the restricted $J$-statistic. This ensures that $C\geq 0$ by construction (as a difference of nested quadratic forms) and that the asymptotic independence result holds.

\bigskip
\begin{remark}[Conditional moments $\Rightarrow$ Unconditional moments]
In the question, the Euler-equation restrictions are stated \emph{conditionally} on the period-$t$ information set:
\[
\E\!\left[m_{t+1}(\theta_0)\mid \mathcal{F}_t\right]=0,
\qquad
m_{t+1}(\theta)=\begin{pmatrix} h_1(\cdot;\theta)R_{t+1}-1 \\ h_2(\cdot;\theta)R_{t+1}-1 \end{pmatrix}.
\]
To obtain the \emph{unconditional} moment restrictions used in GMM, take any instrument vector $z_t$ that is measurable with respect to $\mathcal{F}_t$ (so $z_t\in\mathcal{F}_t$) and apply the law of iterated expectations:
\[
0
=\E\!\left[\E\!\left[m_{t+1}(\theta_0)\mid \mathcal{F}_t\right]\otimes z_t\right]
=\E\!\left[m_{t+1}(\theta_0)\otimes z_t\right].
\]
This a common trick in GMM: multiplying the conditional restriction by an $\mathcal{F}_t$-measurable instrument and then taking expectations converts it into an unconditional moment condition.
\end{remark}

\begin{remark}[Inference with the GMM $J$-test]
The Hansen $J$-test is a \emph{joint} test of the overidentifying restrictions. Under the null hypothesis that \emph{all} moment conditions are correctly specified (and instruments are valid),
\[
H_0:\ \E[g(w_t;\theta_0)]=0,
\]
the (efficient) $J$-statistic
\[
J = n\,\bar g_n(\hat\theta)'\hat S^{-1}\bar g_n(\hat\theta)
\]
has the limiting distribution
\[
J \xrightarrow{d} \chi^2(m-p),
\]
where $m$ is the number of moment conditions and $p$ is the number of parameters. Inference is therefore: reject $H_0$ at level $\alpha$ if $J>\chi^2_{m-p,\,1-\alpha}$ (or report the corresponding $p$-value). A rejection means \emph{at least one} restriction fails (model misspecification and/or invalid instruments), but the $J$-test does not identify \emph{which} moment(s) drive the failure.
\end{remark}



\newpage
%=============================================================================
% SOLUTION TO PROBLEM 2
%=============================================================================
\begin{solution}[GMM objective and test statistics]
Suppose you omit the $\frac{1}{2}$ factor in the usual definition of the GMM objective function and instead use
\[
\tilde{Q}_n(\theta) = -\bar{g}_n(\theta)'\hat{W}\bar{g}_n(\theta),
\]
where $\bar{g}_n(\theta) = \frac{1}{n}\sum_{t=1}^n g(x_t;\theta)$. Suppose you have an optimal weighting matrix (i.e.\ $\hat{W}\xrightarrow{p} S^{-1}$):
\begin{enumerate}
    \item State the asymptotic distribution of the Wald, LM, and QLR tests of $H_0\colon a(\theta_0)=0$ that you would obtain using this modified objective $\tilde{Q}_n$.
    \item Explain how the conclusions of your tests could be affected by omitting the $\frac{1}{2}$ factor. You may assume that all regularity conditions hold.
\end{enumerate}

\medskip
\noindent\textbf{Solution.} Since $\tilde{Q}_n(\theta) = 2\,Q_n(\theta)$ (where $Q_n(\theta)=-\frac{1}{2}\bar{g}_n(\theta)'\hat{W}\bar{g}_n(\theta)$ is the standard objective), the two objective functions differ by a positive multiplicative constant. The key observation is that \textbf{the GMM estimator is unchanged}: both $\hat{\theta}$ (unrestricted) and $\tilde{\theta}$ (restricted by $a(\theta)=0$) are invariant to positive rescaling of the objective, since $\argmin_\theta\,\tilde{Q}_n(\theta)=\argmin_\theta\,Q_n(\theta)$. However, the test statistics that depend on the \emph{value} or the \emph{derivative} of the objective are affected.

\bigskip
\paragraph{(a) Asymptotic distributions of the three tests.} Let us examine each test statistic in turn. Denote $G=\E[\partial g/\partial\theta']$, $S=\Var(\sqrt{n}\bar{g}_n(\theta_0))$, and $\Omega=G'S^{-1}G$.

\medskip
\noindent\textbf{Wald test.} The Wald statistic is
\[
\xi_W = n\,a(\hat\theta)'\bigl[\nabla a(\hat\theta)\,\hat\Omega^{-1}\nabla a(\hat\theta)'\bigr]^{-1}a(\hat\theta),
\]
where $\hat\Omega=\hat G'\hat S^{-1}\hat G$. This statistic depends only on the estimator $\hat\theta$ and its asymptotic variance---it never involves the objective function itself. Since $\hat\theta$ is invariant to the rescaling, the Wald statistic is identical whether we use $Q_n$ or $\tilde Q_n$:
\[
\boxed{\tilde\xi_W = \xi_W \xrightarrow{d} \chi^2(r) \quad\text{under }H_0.}
\]
The Wald test is completely unaffected by the missing $\frac{1}{2}$.

\medskip
\noindent\textbf{LM (Lagrange Multiplier) test.} The LM statistic uses the score of the objective evaluated at the restricted estimator:
\[
\xi_{LM} = n\left(\frac{\partial Q_n(\tilde\theta)}{\partial\theta}\right)'\Omega^{-1}\left(\frac{\partial Q_n(\tilde\theta)}{\partial\theta}\right).
\]
Since $\tilde Q_n=2Q_n$, the gradient scales by 2: $\partial\tilde Q_n/\partial\theta = 2\,\partial Q_n/\partial\theta$. Meanwhile $\tilde\Omega=\Omega$ (it depends on $G$ and $S$, not on the objective scaling). Therefore:
\[
\tilde\xi_{LM} = n\left(\frac{\partial\tilde Q_n(\tilde\theta)}{\partial\theta}\right)'\Omega^{-1}\left(\frac{\partial\tilde Q_n(\tilde\theta)}{\partial\theta}\right) = 4\,\xi_{LM}.
\]
\[
\boxed{\tilde\xi_{LM} = 4\,\xi_{LM} \xrightarrow{d} 4\,\chi^2(r) \quad\text{under }H_0.}
\]

\medskip
\noindent\textbf{QLR (Quasi-Likelihood Ratio) test.} The standard QLR statistic is $\xi_{QLR}=2n[Q_n(\hat\theta)-Q_n(\tilde\theta)]$. With the modified objective:
\[
\tilde\xi_{QLR} = 2n[\tilde Q_n(\hat\theta)-\tilde Q_n(\tilde\theta)] = 2n\cdot 2[Q_n(\hat\theta)-Q_n(\tilde\theta)] = 2\,\xi_{QLR}.
\]
\[
\boxed{\tilde\xi_{QLR} = 2\,\xi_{QLR} \xrightarrow{d} 2\,\chi^2(r) \quad\text{under }H_0.}
\]

\bigskip
\paragraph{(b) Consequences for inference.} If one naively uses the standard $\chi^2(r)$ critical values with the modified test statistics, the LM and QLR tests will \textbf{over-reject} the null hypothesis. The inflated test statistics exceed the critical value more often than they should, leading to excessively high Type~I error rates.

To see this concretely for the QLR test: under $H_0$, the probability of not rejecting is
\[
\Pr\!\left(\tilde\xi_{QLR}<\chi^2_{r,\alpha}\right) = \Pr\!\left(2\,\xi_{QLR}<\chi^2_{r,\alpha}\right) < \Pr\!\left(\xi_{QLR}<\chi^2_{r,\alpha}\right),
\]
so the actual rejection probability exceeds $\alpha$. The same logic applies to the LM test (with an even larger distortion, since the scaling factor is 4 rather than 2).

\medskip
\noindent\textbf{Practical fix.} One can simply divide the LM statistic by 4 and the QLR statistic by 2 to recover the correct asymptotic distributions, or equivalently use critical values from $4\,\chi^2(r)$ and $2\,\chi^2(r)$, respectively. The Wald test requires no adjustment.

\medskip
\noindent\textbf{Why the tests differ.} The Wald statistic is a function of the estimator alone---it measures how far $a(\hat\theta)$ is from zero relative to its estimated variance, and is therefore immune to how the objective is scaled. The LM statistic is a function of the \emph{score} (first derivative) of the objective, so it inherits the factor of 2 in the gradient (and its square inherits the factor of 4). The QLR statistic is a function of the \emph{difference in objective values}, so it inherits the factor of 2 directly. The lesson is straightforward: the $\frac{1}{2}$ in the standard GMM objective is not merely a convention---it is chosen so that the QLR and LM statistics have the same $\chi^2(r)$ limiting distribution as the Wald, preserving their asymptotic equivalence under $H_0$.
\end{solution}



\newpage
%=============================================================================
% SOLUTION TO PROBLEM 3
%=============================================================================
\begin{solution}[Misspecified GMM] Suppose you are estimating an over-identified model by GMM, but the model is misspecified: there is no $\theta\in\Theta$ such that $\E[g(x_t;\theta)]=0$. 
\begin{enumerate}
    \item What is the pseudo-true parameter that you are in fact estimating? Show how it depends on your choice of weighting matrix.
    \item You may assume the pseudo-true parameter is unique and that your estimator is consistent for it. Mimic the proof of asymptotic normality for GMM. Where does the argument break down?
    \item Now suppose the model is ``locally'' misspecified: the degree of misspecification is of order sampling error, $\E[g(x_t;\theta_0)]=c/\sqrt{n}$, where $\theta_0$ is the pseudo-true parameter and $c$ is a vector of constants. Derive the asymptotic distribution of your estimator. How does it differ from the correctly specified case?
    \item Are the usual asymptotic variance estimators valid in this case? Explain.
\end{enumerate}

\medskip
\noindent\textbf{Solution.} Let $\bar{g}_n(\theta)=\frac{1}{n}\sum_{t=1}^n g(x_t;\theta)$ and let $\hat{W}\xrightarrow{p}W$ be a positive definite weighting matrix.

\bigskip
\paragraph{(a) The pseudo-true parameter.} The GMM estimator maximizes (equivalently, minimizes in absolute value) the sample criterion
\[
Q_n(\theta) = -\tfrac{1}{2}\bar{g}_n(\theta)'\hat{W}\bar{g}_n(\theta).
\]
By a uniform law of large numbers, $Q_n(\theta)\xrightarrow{p}Q(\theta)$ where the population criterion is
\[
Q(\theta) = -\tfrac{1}{2}\E[g(x_t;\theta)]'W\,\E[g(x_t;\theta)].
\]
The \textbf{pseudo-true parameter} $\theta_0$ is defined as the maximizer of $Q(\theta)$ over $\Theta$:
\[
\theta_0 = \argmax_{\theta\in\Theta}\; Q(\theta).
\]
The corresponding first-order condition is
\[
0 = \frac{\partial Q(\theta_0)}{\partial\theta} = -G(\theta_0)'W\,\E[g(x_t;\theta_0)],
\]
where $G(\theta)=\E[\partial g(x_t;\theta)/\partial\theta']$ is the $m\times p$ Jacobian. Writing $\bar{g}(\theta)=\E[g(x_t;\theta)]$ for short, this reads
\[
\boxed{G(\theta_0)'W\,\bar{g}(\theta_0) = 0.}
\]

Under correct specification, $\bar{g}(\theta_0)=0$, so this FOC is satisfied trivially and $\theta_0$ does not depend on $W$. Under misspecification, $\bar{g}(\theta_0)=c$ for some $c\neq 0$. The FOC becomes $G(\theta_0)'Wc=0$, and different choices of $W$ will lead to different solutions. \textbf{The pseudo-true parameter depends on the weighting matrix.} Intuitively, $\theta_0$ is the parameter value that makes the moment conditions ``as close to zero as possible'' in the metric defined by $W$---and different metrics yield different notions of ``closest.''

\bigskip
\paragraph{(b) Where asymptotic normality breaks down.} Assume $\theta_0$ is unique and $\hat\theta\xrightarrow{p}\theta_0$. To derive the limiting distribution, we mimic the standard proof. The sample FOC is $\partial Q_n(\hat\theta)/\partial\theta\approx 0$, which gives
\[
0 \approx -G_n(\hat\theta)'\hat{W}\bar{g}_n(\hat\theta).
\]
Taking a mean-value expansion of $\bar{g}_n(\hat\theta)$ around $\theta_0$:
\[
\bar{g}_n(\hat\theta) = \bar{g}_n(\theta_0) + G_n(\tilde\theta)(\hat\theta-\theta_0),
\]
where $\tilde\theta$ lies between $\hat\theta$ and $\theta_0$ (entry-by-entry). Substituting:
\[
G_n(\hat\theta)'\hat{W}\,G_n(\tilde\theta)\sqrt{n}(\hat\theta-\theta_0) = -G_n(\hat\theta)'\hat{W}\sqrt{n}\,\bar{g}_n(\theta_0) + o_p(1).
\]
Under the usual regularity conditions (consistency, interior $\theta_0$, differentiability, dominated Hessian), the left-hand side converges to $(G'WG)\sqrt{n}(\hat\theta-\theta_0)$ plus negligible terms. The question is: \emph{what happens to $\sqrt{n}\,\bar{g}_n(\theta_0)$?}

Under correct specification, $\E[\bar{g}_n(\theta_0)]=0$ and the CLT gives $\sqrt{n}\,\bar{g}_n(\theta_0)\xrightarrow{d}N(0,S)$. Under misspecification, $\E[\bar{g}_n(\theta_0)]=\bar{g}(\theta_0)=c\neq 0$, so
\[
\sqrt{n}\,\E[\bar{g}_n(\theta_0)] = \sqrt{n}\,c \to \infty.
\]
\textbf{The argument breaks down because the leading term on the right-hand side diverges.} The CLT still applies to the centered quantity $\sqrt{n}[\bar{g}_n(\theta_0)-c]$, but the uncentered version $\sqrt{n}\,\bar{g}_n(\theta_0)$ is dominated by the deterministic drift $\sqrt{n}\,c$, which grows without bound. One cannot factor out a well-behaved normal limit and invert $(G'WG)$ to get a standard asymptotic distribution.

\bigskip
\paragraph{(c) Local misspecification.} Now suppose the misspecification is ``small''---of order sampling error:
\[
\E[g(x_t;\theta_0)] = \frac{c}{\sqrt{n}},
\]
where $c$ is a fixed vector of constants and $\theta_0$ is the pseudo-true parameter. This local misspecification framework is attractive because it nests correct specification ($c=0$) and allows us to study what happens when the model is ``almost right'' in a precise asymptotic sense. Returning to the expansion from part (b), we now analyze $\sqrt{n}\,\bar{g}_n(\theta_0)$. Decompose it as
\[
\sqrt{n}\,\bar{g}_n(\theta_0) = \sqrt{n}\bigl[\bar{g}_n(\theta_0)-\E[\bar{g}_n(\theta_0)]\bigr] + \sqrt{n}\,\E[\bar{g}_n(\theta_0)].
\]
The first term converges by the CLT: $\sqrt{n}[\bar{g}_n(\theta_0)-\E[\bar{g}_n(\theta_0)]]\xrightarrow{d}N(0,S)$. The second term is deterministic: $\sqrt{n}\,\E[\bar{g}_n(\theta_0)]=\sqrt{n}\cdot c/\sqrt{n}=c$. Therefore:
\[
\sqrt{n}\,\bar{g}_n(\theta_0) \xrightarrow{d} N(c,\;S).
\]
The crucial difference from part (b) is that the drift $c$ is now a \emph{finite constant} rather than a divergent sequence---the misspecification is exactly calibrated to produce a non-zero but bounded shift in the limit distribution.

Substituting into the expansion and applying Slutsky's theorem:
\begin{align*}
\sqrt{n}(\hat\theta-\theta_0) &= -(G'WG)^{-1}G'W\sqrt{n}\,\bar{g}_n(\theta_0) + o_p(1) \\
&\xrightarrow{d} -(G'WG)^{-1}G'W \cdot N(c,S).
\end{align*}
Since a linear transformation of a normal random vector is normal:
\[
\boxed{\sqrt{n}(\hat\theta-\theta_0) \xrightarrow{d} N\!\left((G'WG)^{-1}G'Wc,\;\;(G'WG)^{-1}G'WSWG(G'WG)^{-1}\right).}
\]

The \textbf{variance} is identical to the correctly specified case (the usual GMM sandwich). The \textbf{mean} is now non-zero: there is an asymptotic bias of $(G'WG)^{-1}G'Wc$, which depends on both the degree of misspecification $c$ and the weighting matrix $W$. Under correct specification ($c=0$), the bias vanishes and we recover the standard result. In essence, local misspecification shifts the center of the limiting distribution without changing its spread.

\bigskip
\paragraph{(d) Validity of the usual variance estimators.} The standard variance estimator is
\[
\hat{S} = \frac{1}{n}\sum_{t=1}^n g(x_t;\hat\theta)\,g(x_t;\hat\theta)'.
\]
Under local misspecification, $\E[g(x_t;\theta_0)]=c/\sqrt{n}$. To check consistency of $\hat S$, we decompose the population second moment using the variance--bias identity:
\begin{align*}
\E\!\bigl[g(x_t;\theta_0)\,g(x_t;\theta_0)'\bigr]
&= \Var\!\bigl(g(x_t;\theta_0)\bigr) + \E\!\bigl[g(x_t;\theta_0)\bigr]\E\!\bigl[g(x_t;\theta_0)\bigr]' \\[4pt]
&= S + \frac{c}{\sqrt{n}}\cdot\frac{c'}{\sqrt{n}} \\[4pt]
&= S + \frac{cc'}{n}.
\end{align*}
The first line is simply $\E[XX'] = \Var(X)+\E[X]\E[X]'$ applied entry-by-entry. In the second line we substituted $\Var(g)=S$ (the long-run variance, which does not depend on $n$) and $\E[g]=c/\sqrt{n}$. Therefore:
\[
\E\bigl[g(x_t;\theta_0)g(x_t;\theta_0)'\bigr] = S + \frac{cc'}{n} \;\xrightarrow{n\to\infty}\; S.
\]
The squared-bias term $cc'/n\to 0$, so the population second moment converges to $S$ as $n\to\infty$. By a standard LLN argument, $\hat{S}\xrightarrow{p}S$, and the \textbf{usual variance estimator remains consistent}. In finite samples, however, $\E[g(x_t;\theta_0)]\neq 0$, so the sample second moment includes a (small) squared-bias component. One can improve finite-sample performance by using a demeaned estimator:
\[
\hat{S}_{\text{dm}} = \frac{1}{n}\sum_{t=1}^n \bigl[g(x_t;\hat\theta)-\bar{g}_n(\hat\theta)\bigr]\bigl[g(x_t;\hat\theta)-\bar{g}_n(\hat\theta)\bigr]',
\]
which subtracts the sample mean of the moments. Since $\bar{g}_n(\hat\theta)=O_p(n^{-1/2})$ under local misspecification, both estimators are asymptotically equivalent, but $\hat{S}_{\text{dm}}$ may be more accurate at moderate sample sizes.

\medskip
\noindent\textbf{Global vs.\ local misspecification: a comparison.} Under \emph{global} misspecification---meaning $\E[g(x_t;\theta_0)]=c\neq 0$ with $c$ fixed and not shrinking---the situation is far worse. The estimator $\sqrt{n}(\hat\theta-\theta_0)$ does not converge to a normal distribution (part~(b)), and the usual variance estimator~$\hat{S}$ is inconsistent for $\Var(g)$ because the squared-bias component $cc'$ does not vanish.

The local misspecification framework avoids these pathologies by making the bias small enough to be ``absorbed'' by the CLT. It is the same asymptotic device used in the weak instruments literature \citep{StaigerStock1997}, where the first-stage coefficient is local-to-zero: $\pi=C/\sqrt{n}$.

\end{solution}

\begin{remark}[Correct specification, global identification, and local identification]
    Throughout this problem set the \textbf{setting} is as follows. We observe an i.i.d.\ (or stationary ergodic) sequence $\{x_t\}$ taking values in $\mathbb{R}^{d_x}$. We have a \emph{moment function}
    \[
    g:\mathbb{R}^{d_x}\times\Theta\;\longrightarrow\;\mathbb{R}^m,
    \qquad \Theta\subseteq\mathbb{R}^p,
    \]
    where $m\geq p$ (we need at least as many moment conditions as parameters). The GMM estimator exploits the \emph{sample analog} of the population moment condition $\E[g(x_t;\theta_0)]=0$.
    
    \medskip\noindent\textbf{Correct specification.} The model is \emph{correctly specified} if there exists $\theta_0\in\Theta$ such that
    \[
    \E[g(x_t;\theta_0)] = 0.
    \]
    When this fails for every $\theta\in\Theta$ the model is \emph{misspecified} (see Problem~3).
    
    \medskip\noindent\textbf{Global identification.} Under correct specification, $\theta_0$ is \emph{globally identified} if it is the \emph{unique} solution to the population moment condition:
    \[
    \E[g(x_t;\theta)]=0 \quad\Longleftrightarrow\quad \theta=\theta_0,
    \qquad\text{i.e.}\quad \E[g(x_t;\theta)]\neq 0\;\;\forall\,\theta\in\Theta,\,\theta\neq\theta_0.
    \]
    Global identification guarantees that the population criterion $Q(\theta)=-\tfrac{1}{2}\E[g(x_t;\theta)]'W\,\E[g(x_t;\theta)]$ is uniquely maximized at $\theta_0$. It is generally \emph{difficult to verify} in nonlinear models because it requires checking all $\theta\in\Theta$.
    
    \medskip\noindent\textbf{Local identification.} $\theta_0$ is \emph{locally identified} if it is the unique solution in some neighborhood $\mathcal{U}$ of $\theta_0$. A sufficient condition (first-order) comes from linearizing the moment function around $\theta_0$:
    \[
    \E[g(x_t;\theta)] \approx \E[g(x_t;\theta_0)] + G(\theta_0)(\theta-\theta_0) = G(\theta_0)(\theta-\theta_0),
    \]
    where $G(\theta_0)=\E\!\left[\tfrac{\partial g(x_t;\theta_0)}{\partial\theta'}\right]\in\mathbb{R}^{m\times p}$ is the Jacobian. For $G(\theta_0)(\theta-\theta_0)=0$ to imply $\theta=\theta_0$, the null space of $G(\theta_0)$ must be trivial, i.e.\
    \[
    \operatorname{rank}(G(\theta_0)) = p.
    \]
    This is the \textbf{rank condition} for local identification (see Problem~4a). It requires $m\geq p$ (order condition) \emph{and} that the $m$ moments collectively provide $p$ linearly independent directions of information about $\theta$.
    
    \medskip\noindent\textbf{How the three concepts relate.}
    \begin{itemize}[leftmargin=1.8em]
      \item Correct specification is a prerequisite: without it, $\theta_0$ is replaced by a pseudo-true parameter that depends on $W$ (Problem~3a).
      \item Global $\Rightarrow$ local identification, but \emph{not} vice versa. A model can be locally identified at $\theta_0$ yet have multiple global solutions.
      \item Without local identification ($\operatorname{rank}(G)<p$), the GMM criterion is flat in at least one direction near $\theta_0$: the estimator is not $\sqrt{n}$-consistent and standard asymptotic theory breaks down.
      \item GMM identification is \emph{weaker} than maximum-likelihood identification: two distinct parameters $\theta\neq\theta'$ may satisfy $\E[g(x_t;\theta)]=\E[g(x_t;\theta')]=0$ yet generate different distributions for $x_t$, so MLE would distinguish them while GMM would not (``limited information'' estimation).
    \end{itemize}
    \end{remark}

\begin{remark}[Two-step optimally-weighted GMM]
In practice the efficient weighting matrix $W=S^{-1}$ is unknown because $S=\Var(\sqrt{n}\,\bar g_n(\theta_0))$ depends on the unknown $\theta_0$. The standard solution is the \textbf{two-step} procedure.

\medskip\noindent\textbf{Step 1 (first-step estimator).} Maximize the GMM criterion using the identity weighting matrix $\hat W = I$:
\[
\tilde\theta \;=\; \argmin_{\theta\in\Theta}\;\bar g_n(\theta)'\,I\,\bar g_n(\theta).
\]
$\tilde\theta$ is consistent for $\theta_0$ (as long as the model is identified), but not efficient.

\medskip\noindent\textbf{Step 2 (optimal weight, second-step estimator).} Use $\tilde\theta$ to construct a consistent estimate of $S$:
\[
\hat W \;=\; \left(\frac{1}{n}\sum_{t=1}^n g(x_t;\tilde\theta)\,g(x_t;\tilde\theta)'\right)^{-1}.
\]
Then re-estimate by maximizing the criterion with this data-dependent weight:
\[
\hat\theta \;=\; \argmin_{\theta\in\Theta}\;\bar g_n(\theta)'\,\hat W\,\bar g_n(\theta).
\]
Because $\hat W\xrightarrow{p}S^{-1}$, the two-step estimator $\hat\theta$ is \textbf{asymptotically efficient} among all GMM estimators that use the same moment conditions, with asymptotic variance $(G'S^{-1}G)^{-1}$.

\medskip\noindent\textbf{Why not iterate?} One can continue re-estimating $\hat W$ with the updated $\hat\theta$ (the \emph{continuously-updated GMM}, CU-GMM, of Hansen, Heaton, and Yaron 1996). CU-GMM shares the same first-order asymptotic distribution as two-step GMM but can have better higher-order properties in some settings.
\end{remark}


\newpage
%=============================================================================
% SOLUTION TO PROBLEM 4
%=============================================================================
\begin{solution}[GMM: identification, efficient weighting, and optimal instruments]
Consider estimation of a parameter vector $\beta\in\R^k$ by GMM applied to the conditional moment condition $\E[e(z_i;\beta_0)\mid x_i]=0$, where $e(z_i;\beta)\in\R$ is continuous and differentiable in $\beta$. Suppose $u_i=(x_i',z_i')'$ is i.i.d. The conditional moment implies the unconditional moment $\E[e(z_i;\beta_0)A(x_i)]=0$ for any vector function $A(x_i)\in\R^m$ of $x_i$. Define $g(u_i;\beta)=e(z_i;\beta)A(x_i)$. The GMM estimator is $\hat{\beta}=\argmin_\beta\,\bar{g}(\beta)'W_n\bar{g}(\beta)$ with $\bar{g}(\beta)=\frac{1}{n}\sum_{i=1}^n g(u_i;\beta)$, where $W_n$ is an $(m\times m)$ symmetric positive definite matrix. Under standard regularity conditions, $\sqrt{n}(\hat{\beta}-\beta_0)\xrightarrow{d}N(0,\Sigma)$ with
\[
\Sigma = (G'WG)^{-1}G'WVWG(G'WG)^{-1},
\]
where $W=\plim W_n$, $G=\E[\frac{\partial}{\partial\beta}g(u_i;\beta_0)]$ $(m\times k)$, and $V=\E[A(x_i)e(z_i;\beta_0)^2A(x_i)']$ $(m\times m)$. The question asks: 
\begin{enumerate}
    \item a sufficient condition for local identification of $\beta_0$;
    \item the infeasible $W_n$ that yields efficient GMM and a proof that $\Sigma-\Sigma^\#\geq 0$;
    \item the asymptotic covariance $\Sigma^*$ when $A(x_i)^*=D(x_i)/\sigma^2(x_i)$ with $D(x_i)=\E[\partial e/\partial\beta\mid x_i]_{\beta_0}$ and $\sigma^2(x_i)=\E[e(z_i;\beta_0)^2\mid x_i]$;
    \item a proof that $\Sigma-\Sigma^*\geq 0$;
    \item how to compute $A^*(x_i)$ in practice.
\end{enumerate}
\medskip
\noindent\textbf{Solution.}
\paragraph{(a) Local identification.} The population GMM first-order condition is $G'W\,\E[g(u_i;\beta)]=0$. For this to locally pin down $\beta_0$, the Jacobian~$G$ must have full column rank. Since $G$ is $m\times k$ ($m\geq k$ is needed), the sufficient condition is
\[
\boxed{\rank(G)=k.}
\]
This is the \textbf{rank condition} for local identification. It requires $m\geq k$ (the order condition: at least as many moment conditions as parameters) together with the requirement that the $m$ moment conditions provide $k$ linearly independent pieces of information about $\beta$. If $\rank(G)<k$, then some directions in parameter space are locally unidentified: small perturbations of $\beta$ along the null space of $G$ leave the moment conditions nearly unchanged, and the GMM objective becomes flat.

\bigskip
\paragraph{(b) Efficient (infeasible) weighting matrix.} Recall the sandwich formula
\[
\Sigma=(G'WG)^{-1}G'WVWG(G'WG)^{-1}.
\]
We look for the $W$ that minimizes $\Sigma$ in the positive-semidefinite ordering.

\medskip
\noindent\textbf{Claim:} The efficient weighting matrix is $W=V^{-1}$, yielding $\Sigma^\#=(G'V^{-1}G)^{-1}$.

\medskip
\noindent\textbf{Verification.} Setting $W=V^{-1}$ in the sandwich formula:
\begin{align*}
\Sigma^\# &= (G'V^{-1}G)^{-1}G'V^{-1}\underbrace{V}_{=V}V^{-1}G(G'V^{-1}G)^{-1} \\
&= (G'V^{-1}G)^{-1}\underbrace{G'V^{-1}G(G'V^{-1}G)^{-1}}_{=I_k} \\
&= (G'V^{-1}G)^{-1}.
\end{align*}
The sandwich collapses because the ``bread'' and ``meat'' simplify. This is a manifestation of the general principle that efficient GMM uses the inverse of the moment covariance as the weighting matrix.

\medskip
\noindent\textbf{Proof of optimality: $\Sigma-\Sigma^\#\geq 0$ for any $W$.} Define
\[
M = (G'WG)^{-1}G'W - (G'V^{-1}G)^{-1}G'V^{-1}.
\]
Note that $MG=I_k-I_k=0$ (both terms, when post-multiplied by $G$, give $I_k$). Now write:
\begin{align*}
\Sigma - \Sigma^\# &= (G'WG)^{-1}G'WVWG(G'WG)^{-1} - (G'V^{-1}G)^{-1}.
\end{align*}
Adding and subtracting cross terms and using $MG=0$, one can verify that
\[
\Sigma - \Sigma^\# = M\,V\,M'.
\]
Since $V$ is positive definite (it is a second-moment matrix of the non-degenerate random vector $e\cdot A$), $M\,V\,M'\geq 0$ in the positive-semidefinite sense. Therefore:
\[
\boxed{\Sigma - \Sigma^\# = M\,V\,M' \geq 0 \quad\text{for any }W.}
\]
Any choice of weighting matrix other than $V^{-1}$ incurs ``excess variance'' $M\,V\,M'$. This weighting matrix is infeasible because $V=\E[A(x_i)e(z_i;\beta_0)^2 A(x_i)']$ depends on the unknown $\beta_0$. In practice, one uses a two-step procedure: first estimate $\hat\beta_1$ with an arbitrary $W$ (e.g., $W=I_m$), then construct $\hat{V}=\frac{1}{n}\sum_i A(x_i)e(z_i;\hat\beta_1)^2 A(x_i)'$, and re-estimate with $\hat{W}=\hat{V}^{-1}$.

\bigskip
\paragraph{(c) Asymptotic covariance with optimal instruments.} The choice of instruments $A(x_i)$ determines $G$ and $V$, and hence the efficiency bound. The claim is that the choice
\[
A^*(x_i) = \frac{D(x_i)}{\sigma^2(x_i)}, \qquad D(x_i)=\E\!\left[\frac{\partial e(z_i;\beta)}{\partial\beta}\,\bigg|\,x_i\right]_{\!\beta=\beta_0}\!\!\!\!\!\!\!, \qquad \sigma^2(x_i)=\E[e(z_i;\beta_0)^2\mid x_i],
\]
yields the smallest possible asymptotic variance among all choices of $A$.

\medskip
\noindent\textbf{Computing $G^*$ and $V^*$.} With $A^*(x_i)=D(x_i)/\sigma^2(x_i)$, define $g^*(u_i;\beta)=e(z_i;\beta)A^*(x_i)$. The Jacobian is:
\begin{align*}
G^* &= \E\!\left[\frac{\partial g^*(u_i;\beta_0)}{\partial\beta'}\right] = \E\!\left[A^*(x_i)\frac{\partial e(z_i;\beta_0)}{\partial\beta'}\right].
\end{align*}
Using the law of iterated expectations (condition on $x_i$ first):
\begin{align*}
G^* &= \E\!\left[\frac{D(x_i)}{\sigma^2(x_i)}\,\E\!\left[\frac{\partial e}{\partial\beta'}\,\bigg|\,x_i\right]\right] = \E\!\left[\frac{D(x_i)\,D(x_i)'}{\sigma^2(x_i)}\right].
\end{align*}
The moment covariance matrix is:
\begin{align*}
V^* &= \E\!\left[A^*(x_i)\,e(z_i;\beta_0)^2\,A^*(x_i)'\right] \\
&= \E\!\left[\frac{D(x_i)}{\sigma^2(x_i)}\,e(z_i;\beta_0)^2\,\frac{D(x_i)'}{\sigma^2(x_i)}\right] \\
&\qquad\qquad\text{(note: $1/\sigma^2$ appears from each $A^*$, giving $1/\sigma^4$ total)}\\
&= \E\!\left[\frac{D(x_i)D(x_i)'}{\sigma^4(x_i)}\,\E[e(z_i;\beta_0)^2\mid x_i]\right] \\
&= \E\!\left[\frac{D(x_i)D(x_i)'}{\sigma^4(x_i)}\cdot\sigma^2(x_i)\right] = \E\!\left[\frac{D(x_i)D(x_i)'}{\sigma^2(x_i)}\right].
\end{align*}
So $G^*=V^*$. Under efficient weighting ($W=V^{*-1}$), the asymptotic variance is:
\[
\Sigma^* = (G^{*\prime}V^{*-1}G^*)^{-1} = (G^{*\prime}G^{*-1}G^*)^{-1} = G^{*-1} = \left(\E\!\left[\frac{D(x_i)D(x_i)'}{\sigma^2(x_i)}\right]\right)^{-1}.
\]
\[
\boxed{\Sigma^* = \left(\E\!\left[\frac{D(x_i)D(x_i)'}{\sigma^2(x_i)}\right]\right)^{-1}.}
\]
This is the \textbf{semiparametric efficiency bound} for estimating $\beta_0$ from the conditional moment restriction $\E[e\mid x]=0$ (Chamberlain, 1987). It is the analog of the Cramér--Rao bound: no regular estimator that uses only the conditional moment restriction can have a smaller asymptotic variance.

\bigskip
\paragraph{(d) Optimality: $\Sigma-\Sigma^*\geq 0$ for any $A$.} We need to show that for \emph{any} instrument function $A(x_i)$, even after using the efficient weighting matrix $W=V^{-1}$ for that particular $A$, the resulting asymptotic variance $\Sigma=(G'V^{-1}G)^{-1}$ is no smaller than $\Sigma^*$.

\medskip
\noindent\textbf{Step 1: Express $G$ and $V$ for arbitrary $A$.} By the law of iterated expectations:
\begin{align*}
G &= \E\!\left[A(x_i)\frac{\partial e(z_i;\beta_0)}{\partial\beta'}\right] = \E\!\left[A(x_i)\,D(x_i)'\right], \\
V &= \E\!\left[A(x_i)\,e(z_i;\beta_0)^2\,A(x_i)'\right] = \E\!\left[\sigma^2(x_i)\,A(x_i)\,A(x_i)'\right].
\end{align*}

\noindent\textbf{Step 2: Apply a matrix Cauchy--Schwarz inequality.} Define the $\R^{m+k}$-valued random vector
\[
\xi_i = \begin{pmatrix} \sigma(x_i)\,A(x_i) \\ D(x_i)/\sigma(x_i) \end{pmatrix}.
\]
The second-moment matrix of $\xi_i$ is positive semidefinite:
\[
\E[\xi_i\xi_i'] = \begin{pmatrix} V & G \\ G' & G^* \end{pmatrix} \geq 0,
\]
where $V=\E[\sigma^2 AA']$, $G=\E[AD']$, and $G^*=\E[DD'/\sigma^2]=V^*$. By the Schur complement condition for positive semidefiniteness:
\[
G^* - G'V^{-1}G \geq 0.
\]
\begin{remark}[Schur complement and positive semidefiniteness]
Let $M$ be a block matrix partitioned as
\[
M = \begin{pmatrix} A & B \\ B' & C \end{pmatrix},
\quad A\in\R^{m\times m},\; C\in\R^{k\times k},\; B\in\R^{m\times k}.
\]
If $A>0$ (positive definite), the \textbf{Schur complement} of $A$ in $M$ is
\[
M/A \;:=\; C - B'A^{-1}B.
\]
The key fact is:
\[
M \geq 0 \quad\Longleftrightarrow\quad A \geq 0 \;\text{ and }\; C - B'A^{-1}B \geq 0.
\]
That is, \emph{the whole block matrix is positive semidefinite if and only if its Schur complement is positive semidefinite} (given $A>0$). Applying this to our second-moment matrix with $A=V$, $B=G$, and $C=G^*$:
\[
\E[\xi_i\xi_i'] = \begin{pmatrix} V & G \\ G' & G^* \end{pmatrix} \geq 0
\quad\Longrightarrow\quad
G^* - G'V^{-1}G \;\geq\; 0.
\]
\end{remark}

\noindent Since $\Sigma=(G'V^{-1}G)^{-1}$ (under efficient weighting for $A$) and $\Sigma^*=G^{*-1}$, and since $G^*-G'V^{-1}G\geq 0$ implies $(G'V^{-1}G)^{-1}\geq G^{*-1}$ by matrix inversion (both matrices are positive definite), we obtain:
\[
\boxed{\Sigma - \Sigma^* \geq 0 \quad\text{for any instrument function }A(x_i).}
\]
The inequality is strict whenever $A$ is not proportional to $A^*$, because the Schur complement is strictly positive definite unless the upper-right block ``explains'' all of $G^*$.

\medskip
\noindent\textbf{Intuition.} The optimal instrument $A^*(x_i)=D(x_i)/\sigma^2(x_i)$ concentrates information where it is most useful: observations with a large expected derivative $D(x_i)$ (high sensitivity of the residual to $\beta$) and small conditional variance $\sigma^2(x_i)$ (low noise) receive more weight. Any other choice of instruments either wastes information (under-weighting informative observations) or introduces noise (over-weighting noisy ones). This is a generalized ``signal-to-noise'' logic: $A^*$ is the conditional analog of the efficient score in maximum likelihood.

\bigskip
\paragraph{(e) Computing $A^*(x_i)$ in practice.} The optimal instruments $A^*(x_i)=D(x_i)/\sigma^2(x_i)$ are infeasible because $D(x_i)=\E[\partial e/\partial\beta'\mid x_i]_{\beta_0}$ and $\sigma^2(x_i)=\E[e(z_i;\beta_0)^2\mid x_i]$ both depend on the unknown $\beta_0$ and on conditional expectations. The standard approach, due to Newey (1990), is a three-step procedure:

\medskip
\noindent\textbf{Step 1: Preliminary estimation.} Obtain a $\sqrt{n}$-consistent estimator $\hat\beta_1$ using any set of valid instruments (e.g., $A(x_i)=x_i$) and standard two-step efficient GMM. This provides a consistent starting point.

\medskip
\noindent\textbf{Step 2: Estimate the conditional expectations nonparametrically.} Using the preliminary estimate $\hat\beta_1$, construct:
\begin{itemize}
    \item The residuals $\hat{e}_i = e(z_i;\hat\beta_1)$ and their derivatives $\partial e(z_i;\hat\beta_1)/\partial\beta'$.
    \item Estimate $D(x_i)$ by regressing $\partial e(z_i;\hat\beta_1)/\partial\beta'$ on $x_i$ nonparametrically (e.g., kernel regression, local polynomials, or sieve/series regression with a polynomial or spline basis in $x_i$). Denote the fitted values $\hat{D}(x_i)$.
    \item Estimate $\sigma^2(x_i)$ by regressing $\hat{e}_i^2$ on $x_i$ nonparametrically. Denote the fitted values $\hat\sigma^2(x_i)$.
\end{itemize}

\noindent\textbf{Step 3: Re-estimate with constructed optimal instruments.} Form
\[
\hat{A}^*(x_i)=\frac{\hat{D}(x_i)}{\hat\sigma^2(x_i)}
\]
and use these as instruments in a final round of GMM:
\[
\hat\beta^* = \argmin_\beta\;\bar{g}^*(\beta)'\hat{W}^*\bar{g}^*(\beta), \qquad \bar{g}^*(\beta) = \frac{1}{n}\sum_{i=1}^n e(z_i;\beta)\hat{A}^*(x_i).
\]
Under regularity conditions (including appropriate rates for the nonparametric estimates), Newey (1990) shows that $\hat\beta^*$ achieves the semiparametric efficiency bound $\Sigma^*$ from part (c). The key requirement is that the nonparametric estimates converge fast enough that the estimation error in $\hat{A}^*$ does not affect the first-order asymptotics of $\hat\beta^*$.

\medskip
\noindent\textbf{Practical remarks.}
\begin{itemize}
    \item If the conditional expectation functions $D(\cdot)$ and $\sigma^2(\cdot)$ are believed to be smooth, series regression with polynomial or spline bases works well and is simple to implement.
    \item If the model is linear ($e(z_i;\beta)=y_i-x_i'\beta$), then $D(x_i)=-x_i$ is known and only $\sigma^2(x_i)=\E[(y_i-x_i'\beta_0)^2\mid x_i]$ needs to be estimated. The optimal instruments become $A^*(x_i)=-x_i/\sigma^2(x_i)$, which is the heteroskedasticity-efficient instrument from the GLS literature.
    \item In time-series applications, one must also account for serial correlation in the long-run variance estimation, but the logic of optimal instruments carries over with the appropriate modifications.
\end{itemize}

\end{solution}



\vspace{1.5em}
\begin{thebibliography}{10}
\bibitem[Chamberlain(1987)]{Chamberlain1987}
Chamberlain, G.\ (1987).
\newblock ``Asymptotic efficiency in estimation with conditional moment restrictions.''
\newblock \emph{Journal of Econometrics} 34(3), 305--334.

\bibitem[Eichenbaum et al.(1988)]{EichenbaumHansenSingleton1988}
Eichenbaum, M., L.\ P.\ Hansen, and K.\ J.\ Singleton (1988).
\newblock ``A time series analysis of representative agent models of consumption and leisure choice under uncertainty.''
\newblock \emph{Quarterly Journal of Economics} 103(1), 51--78.

\bibitem[Hansen(1982)]{Hansen1982}
Hansen, L.\ P.\ (1982).
\newblock ``Large sample properties of generalized method of moments estimators.''
\newblock \emph{Econometrica} 50(4), 1029--1054.

\bibitem[Newey(1985)]{Newey1985}
Newey, W.\ K.\ (1985).
\newblock ``Generalized method of moments specification testing.''
\newblock \emph{Journal of Econometrics} 29(3), 229--256.

\bibitem[Newey(1990)]{Newey1990}
Newey, W.\ K.\ (1990).
\newblock ``Efficient instrumental variables estimation of nonlinear models.''
\newblock \emph{Econometrica} 58(4), 809--837.

\bibitem[Staiger and Stock(1997)]{StaigerStock1997}
Staiger, D.\ and J.\ H.\ Stock (1997).
\newblock ``Instrumental variables regression with weak instruments.''
\newblock \emph{Econometrica} 65(3), 557--586.
\end{thebibliography}

\end{document}
